\bookchapter{Bab 33: Pandangan Hadapan — Perjalanan Seterusnya Telefon}

\chapterintro{Jika telefon kesan jatuh, telefon akan papar amaran, jika tidak respons dalam 60 saat, auto-panggil kecemasan.}

\booksection{AR/VR dan pengkomputeran spatial: “Meja kerja kedua” di luar telefon}

Sesetengah telefon Android (Samsung, Google Pixel) juga ada fungsi pengesanan jatuh.

\begin{itemize}
  \item \textbf{Vision Pro susulan}: Pelancaran Apple Vision Pro menandakan konsep "pengiraan spatial" secara rasmi mendarat, tetapi harga generasi pertama (29999 RMB ke atas) dan berat masih halangan. Dijangka 2026-2027 akan keluarkan versi lebih ringan, lebih murah
  \item \textbf{2. "Halimunan" Telefon dan Asas Identiti Baharu}: Buka "desktop kedua" di depan mata anda, maklumat digital menjadi seperti udara di mana-mana. Kerjasama Meta dan Ray-Ban cermin pintar telah jual berjuta-juta, domestik Rokid, XREAL, INMO dll pengeluar juga keluarkan produk. Cermin mata AR kini boleh papar navigasi HUD, dialog penolong AI, tetapi masih terhad paparan waveguide optik dan jangka hayat bateri, jarak "ganti telefon" masih 3-5 tahun.
  \item \textbf{Jam Android: Fungsi kesihatan dan kecemasan serupa.}: Samsung Galaxy Ring, Oura Ring dan cincin pintar lain mula popular, sebagai "sensor berhampiran badan" telefon menambah pemantauan kesihatan. Peranan telefon sedang dari "terminal serba boleh" ke "hab + berbilang aksesori ringan" kombinasi.
\end{itemize}

\booksection{Antara muka otak-komputer: Fikiran menjadi medium interaksi}

VI. Penyelamatan diri dalam situasi berbahaya

\begin{itemize}
  \item \textbf{1. Serangan}Aktifkan SOS senyap: Tekan butang kuasa 3 kali untuk aktifkan SOS senyap, telefon akan rakam audio/video dan hantar lokasi kepada kontak kecemasan tanpa bunyi.
  \item \textbf{2. Hilang di tempat asing}Guna telefon untuk muat turun peta luar talian, pastikan boleh navigasi tanpa rangkaian.
  \item \textbf{Hantar lokasi real-time kepada keluarga.}: Antara muka otak-komputer jarak interaksi telefon peringkat pengguna masih 10-20 tahun, tetapi jangka pendek mungkin dahulu guna untuk bantuan kebolehcapaian, pemantauan fokus dll senario. Isu etika dan privasi ("gelombang otak saya dianggap privasi siapa") akan lebih awal daripada kematangan teknologi menjadi tumpuan perdebatan.
\end{itemize}

\booksection{“Samaran” telefon dan asas identiti baharu}

Simpan kuasa: Tutup App tidak perlu, kurangkan kecerahan, aktifkan mod jimat kuasa.

\begin{itemize}
  \item \textbf{Daripada "pusat peranti" ke "pusat identiti"}Bab 24: Masa depan telefon pintar
  \item \textbf{3. Pemikiran Etika: Adakah Kita Masih Tuannya Telefon?}Telefon pintar terus berkembang dengan teknologi baharu. Bab ini membincangkan trend dan teknologi yang akan datang.
  \item \textbf{Apabila algoritma lebih memahami apa yang anda ingin lihat berbanding diri sendiri, berapa banyak "suka" anda yang benar-benar spontan?}I. Skrin boleh lipat
\end{itemize}

\booksection{Refleksi etika: Masihkah kita tuan kepada telefon?}

Jadilah pengagih sumber, bukan penerima dorongan.

\begin{itemize}
  \item \textbf{Telefon seperti Samsung Galaxy Z Fold, Huawei Mate X, vivo X Fold, OPPO Find N menggunakan skrin yang boleh dilipat seperti buku.}Selepas dibuka, skrin lebih besar (7.6-8 inci), sesuai untuk produktiviti dan multimedia.
  \item \textbf{2. Skrin lipat clamshell}Telefon seperti Samsung Galaxy Z Flip, Huawei Pocket, OPPO Find N Flip menggunakan skrin yang boleh dilipat seperti clamshell.
  \item \textbf{Selepas dilipat, lebih kecil dan mudah dibawa.}: Telefon AI, skrin lipat, komunikasi satelit semua sedang memperlebar jurang antara "mampu guna" dan "tidak mampu guna". Akses adil infrastruktur digital, akan menjadi isu sosial baharu - seperti 20 tahun lalu "tahu tidak tahu internet" sama penting.
  \item Jadikan automasi "pembantu kopilot" anda yang setia.
  \item Bina "garisan belakang digital" anda yang kukuh.
\end{itemize}

\booksection{Kekal bersambung, kekal berdikari — Falsafah hidup digital anda}

II. Bateri generasi baharu

\begin{enumerate}
  \item Sentiasa jadikan keselamatan dan privasi sebagai prasyarat segala kemudahan.
  \item Mari cuba: Susun "Terma Perkhidmatan Penggunaan Peranti Digital Peribadi" anda
  \item Tulis satu "terma perkhidmatan" ringkas untuk anda dan telefon anda, untuk mentakrifkan hubungan anda. Contohnya:
  \item Perkara 1: Masa perhatian saya milik saya. Anda harus memasuki mod "jangan ganggu" secara proaktif semasa saya bekerja/menemani keluarga.
  \item Perkara 2: Data peribadi saya adalah harta saya. Tanggungjawab utama anda ialah penyulitan keselamatan dan penyegerakan mudah.
\end{enumerate}

\begin{tcbpractice}{Perkara 3: Matlamat bersama kita ialah meningkatkan kecekapan dan pengalaman hidup. Jika anda membuat saya ketagihan melalui reka bentuk menggoda, saya berhak menyahpasang anda.}
1. Bateri silikon-karbon

\begin{itemize}
  \item Baiklah, setelah "perjalanan panjang" lebih dua puluh bab sebelum ini, tahniah! Anda telah berjaya meningkat dari "pendatang baru telefon pintar" dan menguasai kemahiran teras dari asas hingga melangkaui. Lampiran ini adalah seperti buku rujukan pantas dan kotak sumber yang disediakan untuk anda.
  \item Lampiran A: Proses Unboxing Telefon Baru – Panduan "Peringkat Pengasuh"
  \item 1. "Tiga Langkah" Pemeriksaan Unboxing:
\end{itemize}

\end{tcbpractice}

\begin{tcbsummary}{Windows 11: Tetapan Sistem Papar Pelbagai paparan Sambung ke paparan wayarles.}
\begin{itemize}
  \item Vision Pro dan cermin mata AR adalah pintu masuk pengiraan baharu selain telefon, tetapi jangka pendek sukar gantikan telefon
  \item 2. Bateri pepejal
  \item Bateri pepejal menggunakan elektrolit pepejal bukan cecair, menawarkan: ketumpatan tenaga lebih tinggi, keselamatan lebih baik, jangka hayat lebih panjang.
  \item Dijangka mula diguna dalam telefon pintar selepas 2026-2027.
  \item 3. Cas pantas baharu
\end{itemize}

\end{tcbsummary}
